% One-file source. Compile twice with pdflatex.
\documentclass[11pt]{article}
\usepackage[a4paper,margin=27mm,top=25mm,bottom=25mm,headheight=14pt]{geometry}
\usepackage[T1]{fontenc}
\usepackage{lmodern,microtype,amsmath,amssymb,amsthm,mathtools,booktabs,tabularx,fancyhdr,xcolor,hyperref}
\definecolor{deep}{HTML}{165D64}
\hypersetup{colorlinks=true,linkcolor=deep,citecolor=deep,urlcolor=deep,pdftitle={Where the periodic fractional Green function first vanishes}}
\pagestyle{fancy}\fancyhf{}\fancyhead[L]{\small\sffamily Fractional Green functions}\fancyhead[R]{\small\sffamily Research manuscript}\fancyfoot[C]{\small\thepage}
\theoremstyle{plain}\newtheorem{theorem}{Theorem}[section]\newtheorem{lemma}[theorem]{Lemma}\newtheorem{corollary}[theorem]{Corollary}\newtheorem{proposition}[theorem]{Proposition}
\newcommand{\R}{\mathbb R}\newcommand{\T}{\mathbb T}\newcommand{\Z}{\mathbb Z}
\newcommand{\dd}{\,\mathrm d}
\numberwithin{equation}{section}
\setlength{\emergencystretch}{2em}
\begin{document}
\thispagestyle{plain}
\noindent{\small\sffamily\bfseries\color{deep} ANALYSIS \quad / \quad 16 NOVEMBER 2026}\par\vspace{9pt}
\noindent{\sffamily\bfseries\fontsize{26}{31}\selectfont Where the periodic fractional Green function first vanishes}\par\vspace{7pt}
\noindent{\large The Le--Pelinovsky conjectures, stabilization of real zeros, and the limit of fourth order}\par\vspace{10pt}
\noindent{\small Research manuscript prepared for Keio}\par
\begin{abstract}
For $2<\alpha<4$ we establish the zero-count conjectures of Le and Pelinovsky for the multiplier $(c+|\xi|^\alpha)^{-1}$ on the circle and the line. A centered-arc inequality forces every nonnegative periodic Green function to be strictly decreasing on the open half-circle. At the first loss of positivity it has exactly one zero, a quadratic contact at the antipode. At large $c$ its real zero count is eventually constant and counts precisely the simple zeros of the line kernel. We derive a two-term logarithmic asymptotic for this integer as $\alpha\uparrow4$ and a three-term expansion for the first-contact parameter as $\alpha\downarrow2$. The fourth-order case was resolved by Le and Pelinovsky, and the negative algebraic tail on the line was already established by Stefanov and Kevrekidis. An independent appendix gives the sharp Bessel-zero cutoff for positive moment measures of numerical Jacobi--Stirling rows, starting from a counterexample already in the corpus.
\end{abstract}
\section{The original questions}
Let $\T=\R/(2\pi\Z)$ and use Lebesgue measure for convolution. The Green functions in Conjectures 1.1 and 1.2 of Le and Pelinovsky \cite{lp} are
\begin{equation}\label{def}
 G_c(x)=\frac1{2\pi}\sum_{n\in\Z}\frac{e^{inx}}{c+|n|^\alpha},
 \qquad H_c(x)=\frac1{2\pi}\int_{\R}\frac{e^{ix\xi}}{c+|\xi|^\alpha}\dd\xi .
\end{equation}
The exponent is fixed unless explicitly varied. On the circle the authors ask for a positive, decreasing small-$c$ regime, finiteness of the number of zeros in the remaining regime, and boundedness of that count as $c\to\infty$ for $2<\alpha<4$, with unbounded growth for $\alpha=4$. On the line they ask for loss of positivity and monotonicity at every $c>0$, and for finite versus infinite zero counts in the same two exponent ranges. The authors already give the fourth-order solution in their Appendix B. Simon \cite[Remark (e)]{simon} explicitly singles out the outstanding periodic zero count. Stefanov and Kevrekidis \cite[Lemma 6]{sk} obtain the negative tail of the line kernel. That tail furnishes a starting sign and does not count zeros of the periodic sum by itself.
\begin{theorem}\label{main}
For each $2<\alpha<4$, the following statements hold.
\begin{enumerate}
\item There is a unique $c_*(\alpha)>0$ separating strict positivity of $G_c$ for $0<c<c_*$ from sign change for $c>c_*$. For $0<c\le c_*$, $G_c'(x)<0$ on $(0,\pi)$. The only zero at $c_*$ is a double zero at $\pi$. Two simple zeros emerge there as $c$ increases.
\item $H_c$ takes both signs, fails to decrease monotonically on the positive half-line, and has finitely many real zeros for every $c>0$. Its positive zeros are simple crossings or double local maxima.
\item $G_c$ has finitely many real zeros for every $c>0$. For all sufficiently large $c$, their number on one circle is exactly the number of simple real zeros of $H_1$ on $\R$. All these periodic zeros are simple. The resulting integer $N(\alpha)$ satisfies, with $U=\log (1/(4-\alpha))$,
\[
             N(\alpha)=\frac{2}{\pi}\bigl(U+5\log U\bigr)+O(1)
             \qquad(\alpha\uparrow4).
\]
\end{enumerate}
\end{theorem}
\paragraph{Auxiliary problem.} Given an even line kernel with finitely many zeros, which of those zeros survive periodization as the period tends to infinity? A tangency and a crossing have different answers. Section~\ref{sec:images} proves a general criterion using the common sign of all noncentral translates. Section~\ref{sec:arc} treats a second issue that the limit on the line cannot decide: which point first reaches zero on a finite circle? The corpus supplied two relevant proof patterns. Manuscript [026] compares zeros in disks under rescaling, and [175] separates a bounded interval from a tail with a proven sign \cite{vmrc}. Here a strict sign for the image sum and a centered-arc inequality give the additional information.
\section{A signed formula for the line kernel}\label{sec:line}
Put $H=H_1$, $\theta=\pi/\alpha$, $u=\cos\theta$, $v=\sin\theta$, $p=1+\alpha$, and $d=-\sin(\pi\alpha/2)/\pi>0$. For $x>0$,
\begin{align}
 H(x)&=\frac{2}{\alpha}e^{-vx}\sin(ux+\theta)-B(x),\label{split}\\
 B(x)&=d\int_0^\infty e^{-rx}
       \frac{r^\alpha}{1+2r^\alpha\cos(\pi\alpha/2)+r^{2\alpha}}\dd r .\label{B}
\end{align}
To verify this, rotate the ray of $e^{ixz}/(1+z^\alpha)$ counterclockwise from the positive real to the positive imaginary axis using the principal power. The only enclosed pole is $w=e^{i\theta}$; its residue is $-we^{ixw}/\alpha$. The outer arc integral tends to zero since $\alpha>1$ and $|e^{ixz}|\le1$ there. The inner arc about zero tends to zero as well. Take real parts and divide by $\pi$. The imaginary-axis integrand contributes $\Re(i/(1+r^\alpha e^{i\pi\alpha/2}))$, which has the negative sign displayed in \eqref{split}.
The denominator in \eqref{B} has a strictly positive lower bound for fixed $\alpha$. Changing variables $s=rx$ and applying dominated convergence gives
\begin{equation}\label{tail}
 B(x)\sim d\,\Gamma(p)x^{-p},\qquad H(x)\sim-d\,\Gamma(p)x^{-p}.
\end{equation}
The Laplace integral is holomorphic in $\Re x>0$. On strips $|\Im x|\le M$, $\Re x\ge1$, the formula gives $H(x)=O_M((\Re x)^{-p})$. We use the even reflection to extend to the opposite half-plane. Continuity on the real axis and $H(0)>0$ follow directly from \eqref{def}. Thus zeros lie in a compact subinterval of $(0,\infty)$ and its reflection, where analyticity gives finiteness. Eventual negativity and the limit zero imply a failure of monotone decrease. Scaling gives $H_c(x)=c^{1/\alpha-1}H(c^{1/\alpha}x)$.
Manuscript [127] of \cite{vmrc} counts crossings within a single oscillatory lobe for a fractional eigenfunction. For the Fourier resolvent here, multiplication by a real exponential gives the following curvature test.
\begin{lemma}\label{hump}
In each component of $\{x>0:\sin(ux+\theta)>0\}$ the function $e^{vx}H(x)$ is strictly concave with strictly negative second derivative. Every multiple positive zero of $H$ is a double strict local maximum. A complete such component has exactly two simple zeros if its maximum after multiplication by $e^{vx}$ is positive.
\end{lemma}
\begin{proof}
The second derivative of $e^{vx}B(x)$ is the integral in \eqref{B} with exponential $e^{(v-r)x}$ and extra factor $(v-r)^2$, hence is positive. The sine has strictly negative second derivative on its positive components. On their complement $H<0$. The claims follow from strict concavity and the negative endpoint values on a complete component.
\end{proof}
\section{Which zeros survive periodization?}\label{sec:images}
\begin{proposition}\label{images}
Let $h$ be real, continuous, even, positive at zero, and strictly negative on $[R,\infty)$ for some $R>0$. Assume that its extension on the positive half-line is holomorphic in $\Re z>0$, with $h(z)=O_M((\Re z)^{-1-\delta})$ for $\Re z\ge1$, $|\Im z|\le M$ and some $\delta>0$. If every positive zero is simple or a strict local maximum, then for all sufficiently large $L$ the periodization $P_L(x)=\sum_{k\in\Z}h(x+kL)$ has exactly as many zeros per period as $h$ has simple real zeros. Each surviving zero is simple.
\end{proposition}
\begin{proof}
Take $L>2R$. For $R\le x\le L/2$ all translates have negative values. On $[0,R]$ the correction $P_L-h$ is strictly negative and $O(L^{-1-\delta})$ uniformly. The periodization is therefore positive in a fixed neighborhood of zero for sufficiently large $L$. Around each positive root in $(0,R)$ take a disjoint small interval. At a touching maximum the original function is nonpositive throughout that interval and every image correction is strictly negative. No real zero remains there. Around a simple zero one can use a small complex disk within the right half-plane. The analytic correction converges uniformly to zero on the disk by the strip estimate and even reflection of the negative translates. Rouch\'e's theorem gives one zero counted with multiplicity; conjugation symmetry makes it real. The remaining compact set is bounded away from zero in absolute value. Reflection accounts for the negative half-period. The assumptions ensure the original list of zeros is finite.
\end{proof}
Apply Proposition~\ref{images} to $H$ using Lemma~\ref{hump}. Write $L=2\pi c^{1/\alpha}$. Fourier coefficients give the exact identity
\begin{equation}\label{period}
              G_c(x)=c^{1/\alpha-1} P_L(c^{1/\alpha}x).
\end{equation}
Indeed the $n$th Fourier coefficient of $P_L$ is $L^{-1}(1+|2\pi n/L|^\alpha)^{-1}$ by absolute integrability and Fourier inversion. Both sides are continuous. For any fixed $L>0$, $P_L$ is real analytic away from $L\Z$ by normal convergence on compact subsets there with the appropriate reflected continuations. At the lattice point $G_c(0)>0$ by the Fourier series. Thus each fixed period has only finitely many zeros. For large $c$ there is sign change by the negative tail and positivity at zero. There is also a failure of monotone decrease on the half-circle: choose fixed $0<x_1<x_2$ with $H(x_1)<H(x_2)<0$ and use uniform convergence of $P_L$ at these points. This proves the zero-count and large-$c$ profile assertions for $2<\alpha<4$.
\section{First contact on the circle}\label{sec:arc}
For any $\alpha>2$ let $\beta=\alpha-2$, and let $K_\beta$ be the even, zero-mean convolution kernel with Fourier multiplier $|n|^{-\beta}$ at $n\ne0$. It belongs to $L^1(\T)$ and strictly decreases on $(0,\pi)$. Indeed, writing $P_r(x)=(1-r^2)/(1-2r\cos x+r^2)$,
\[
 K_\beta(x)=\frac{1}{2\pi\Gamma(\beta)}\int_0^\infty
       t^{\beta-1}(P_{e^{-t}}(x)-1)\dd t .
\]
The $L^1$ norm of the parenthesis is bounded for small $t$ and $O(e^{-t})$ for large $t$. On compact arcs away from zero the integral converges pointwise to a continuous function; each Poisson kernel strictly decreases there between zero and $\pi$. Its Fourier coefficients give the stated multipliers.
\begin{proposition}\label{arc}
If $G_c\ge0$ on the circle then $G'_c(x)<0$ for every $0<x<\pi$. If in addition $G_c(\pi)=0$, then $G_c''(\pi)>0$ and $\partial_cG_c(\pi)<0$.
\end{proposition}
\begin{proof}
The measure $cG_c(y)\dd y$ has mass one. Fourier multiplication gives the distributional identity
\begin{equation}\label{potential}
                -G_c''=K_\beta-c K_\beta*G_c .
\end{equation}
Here $G_c\in C^1$ by absolute Fourier convergence and the right side belongs to $L^1$. For $0<x<\pi$, let $A_x(y)=\int_{-x}^x K_\beta(z-y)\dd z$. This even function attains its strict maximum at $y=0$: on $(0,\pi)$ its a.e. derivative is $K_\beta(y+x)-K_\beta(y-x)<0$, since the circular distance from zero to $y+x$ exceeds that to $y-x$. It is absolutely continuous, including across the possible integrable singularity. Integrating \eqref{potential} over $[-x,x]$ gives
\[
      -2G_c'(x)=\int_{\T} cG_c(y)\,[A_x(0)-A_x(y)]\dd y>0 .
\]
At $\pi$, $K_\beta$ has a strict minimum. Thus \eqref{potential} gives $G_c''(\pi)>0$ by the same averaging argument, with pointwise evaluation valid by continuity away from zero. Differentiating the Fourier series gives $\partial_c G_c=-G_c*G_c$. If $G_c(\pi)=0$, strict monotonicity and evenness make it positive everywhere else, so the convolution at $\pi$ is strictly positive.
\end{proof}
The kernel is strictly positive for sufficiently small $c$, since the constant Fourier mode dominates all others when $1/c>2\zeta(\alpha)$. Nonnegativity is downward closed in $c$. Indeed, if $G_b\ge0$ and $0<c<b$, its resolvent expansion is
\[
 G_c=\sum_{j\ge0}(b-c)^j G_b^{*(j+1)} .
\]
It converges uniformly because $\|G_b\|_1=1/b$ and $\|G_b^{*(j+1)}\|_\infty\le\|G_b\|_\infty b^{-j}$. The equality follows mode by mode. Since $G_b>0$ near zero, a sufficiently high convolution power is strictly positive everywhere on the circle (the repeated sum of that open arc covers the circle). Thus $G_c>0$ whenever $c<b$. There is an independent finite upper bound valid for all $\alpha>2$. If $G_c\ge0$, integrate $1-\cos(2x)\le4(1-\cos x)$ against the probability density $cG_c$, which is positive near zero. The inequality is strict after integration. Using its Fourier coefficients gives
\[
       \frac{2^\alpha}{c+2^\alpha}<\frac4{c+1},
       \qquad c<\frac{3\,2^\alpha}{2^\alpha-4}.
\] Continuity in $c$ in the uniform norm now gives a positive finite supremum $c_*$ of the nonnegative range. At $c_*$ the kernel is nonnegative and has a zero, since a strictly positive minimum would persist above $c_*$. Proposition~\ref{arc} puts its only zero at $\pi$, and proves quadratic contact. Above $c_*$ the kernel changes sign since $G_c(0)>0$. All arguments in this section hold for every $\alpha>2$. In particular the same first-contact theorem extends beyond fourth order, with
\[
       \frac{1}{2\zeta(\alpha)}\le c_*(\alpha)
           <\frac{3\,2^\alpha}{2^\alpha-4}.
\]
\paragraph{The first two zeros.} Write $g=G_{c_*}$ and $h=c-c_*>0$. The zeros immediately above the threshold are
\[
 x=\pi\ \pm\left(\frac{2(g*g)(\pi)}{g''(\pi)}\right)^{1/2}
             \sqrt h+o(\sqrt h) \pmod{2\pi}.
\]
Indeed \eqref{potential} gives continuous dependence of $G_c''$ on $(c,x)$ near $(c_*,\pi)$. Evenness about $\pi$ makes $G_c'(\pi)=0$ for every $c$. A fixed small neighborhood remains strictly convex; its boundary and the complementary compact arc remain positive. Taylor's formula at $\pi$ and the strictly negative parameter derivative give the asserted roots and uniqueness.
\section{Growth of the limiting count near fourth order}
Let $\varepsilon=4-\alpha\downarrow0$. The integral in \eqref{B} has uniform two-sided estimates: there are $m,M>0$ such that, for $x\ge1$ and $7/2\le\alpha<4$,
\[
          m\varepsilon x^{-p}\le B(x)\le M\varepsilon x^{-p}.
\]
In fact $d\asymp\varepsilon$, the denominator is at least one (the cosine is nonnegative in this range), and the upper bound follows by the gamma integral. For the lower bound integrate over $0<r<1/x$, where the denominator is at most four. All constants can be chosen independently of the exponent.
Let $U=\log(1/\varepsilon)$ and $X=(U+p\log U)/v$. The positive function $D(x)=\varepsilon e^{vx}x^{-p}$ has value bounded above and below away from zero at $X$, uniformly as $U\to\infty$. Its logarithmic derivative $v-p/x$ is bounded below by a fixed positive number for $x\ge x_0$, with $x_0$ absolute. Hence for some fixed $C$, the normalized correction $e^{vx}B(x)$ is less than $2/\alpha$ for $x_0\le x\le X-C$, and exceeds $2/\alpha$ everywhere beyond $X+C$. Every complete positive sine lobe with maximum point in $[x_0,X-C]$ therefore contains exactly two simple crossings by Lemma~\ref{hump}. No zeros occur past $X+C$. The initial interval and the transition window account for $O(1)$ zeros, since each positive lobe has at most two and $u$ stays bounded away from zero. The sine maxima are $2\pi/u$ apart, so the full-line count of simple zeros is
\[
       N(\alpha)=\frac{2u}{\pi v}(U+p\log U)+O(1).
\]
Here $u/v=1+O(\varepsilon)$ and $p=5-\varepsilon$, which gives Theorem~\ref{main}. At $\alpha=4$ the correction in \eqref{B} disappears and $H(x)=\tfrac12e^{-|x|/\sqrt2}\sin(|x|/\sqrt2+\pi/4)$. Hence the line has infinitely many simple zeros. Each fixed finite set of its sign changes persists for all large periods by convergence on compact intervals, so the periodic counts tend to infinity. Finiteness per fixed period follows by analyticity off the lattice and positivity at zero, as before. This also recovers the fourth-order conclusions of \cite[Appendix B]{lp}.
\section{The first-contact parameter as the order decreases to two}
The authors' Appendix C derives a formal leading approximation to the first root of $G_c(\pi)$ and explicitly leaves its error terms uncontrolled. The first-contact theorem identifies that root with $c_*$. A Fourier expansion at the ordinary second-order resolvent gives a sharper result.
\begin{theorem}\label{threshold-asym}
Set $\varepsilon=\alpha-2$ and $V=\log(1/\varepsilon)$. As $\varepsilon\downarrow0$,
\[
 \pi\sqrt{c_*(2+\varepsilon)}=V+3\log V+\log\!\frac{4}{7\zeta(3)}+o(1).
\]
In particular the leading law $c_*\sim V^2/\pi^2$ has the coefficient in \cite[Appendix C, (C.7)]{lp}.
\end{theorem}
\begin{proof}
Write $s=\sqrt c$ and indicate the order by a superscript. At order two the periodic sum of the line kernel $e^{-s|x|}/(2s)$ gives $G^{(2)}_{s^2}(\pi)=1/(2s\sinh(\pi s))$. Uniformly for $s\ge1$, Taylor expansion in the order yields
\begin{align}
 G^{(2+\varepsilon)}_{s^2}(\pi)&=\frac1{2s\sinh(\pi s)}
          -\frac{\varepsilon}{\pi s^4}A(s)+O(\varepsilon^2),\label{expand}\\
 A(s)&=\sum_{n\ge1}(-1)^n n^2\log n\,(1+n^2/s^2)^{-2}.
\end{align}
Indeed, both derivatives through order two may be taken termwise: the absolute value of the second derivative of $(s^2+n^\alpha)^{-1}$ is bounded by $3n^{-2}(\log n)^2$ for $\alpha\ge2$.
  
We show $A(s)\to A_0=7\zeta(3)/(4\pi^2)$. Here is one regularization argument with a summable majorant. With $D f(n)=f(n)-f(n+1)$, put $S(f)=\sum_{n\ge1}(-1)^{n-1}f(n)$. For decaying summable data, five iterations of $2S(f)=f(1)+S(Df)$ give
\begin{equation}\label{altsum}
 S(f)=\sum_{j=0}^4 2^{-j-1}D^jf(1)+2^{-5}S(D^5 f).
\end{equation}
Apply this to $f_s(x)=x^2\log x\,(1+x^2/s^2)^{-2}$. For $x\ge1$ the $k$th derivative of the weight is bounded by $C_k x^{-k}$ uniformly in $s$: the functions $y^k(\mathrm d/\mathrm dy)^k(1+y^2)^{-2}$ are bounded. Leibniz's rule then gives $|f_s^{(5)}(x)|\le Cx^{-3}(1+\log x)$. The fifth difference is a fivefold integral of this derivative over a translated unit cube, up to sign. Thus the right side of \eqref{altsum} converges to the same expression with $f(x)=x^2\log x$. To identify the value, apply that right side to $x^{-z}$: it defines a holomorphic function for $\Re z>-4$ by the fifth-difference estimate, and equals $\eta(z)=(1-2^{1-z})\zeta(z)$ for $\Re z>1$. Differentiation at $z=-2$ and the opposite sign in the definition of $A$ give $A_0=\eta'(-2)=-7\zeta'(-2)$. The zeta functional equation gives $\zeta'(-2)=-\zeta(3)/(4\pi^2)$.
  
For each fixed $c>0$ the Fourier series converges uniformly to the strictly positive second-order kernel as $\alpha\downarrow2$, so $s_*:=\sqrt{c_*(2+\varepsilon)}\to\infty$. Evaluate \eqref{expand} first at $s=2V/\pi$. The positive term is $O(\varepsilon^2/s)$, and the negative term has magnitude asymptotic to $\varepsilon A_0/(\pi s^4)$, dominating both errors. Thus $s_*\le 2V/\pi$. In particular $\varepsilon s_*^4\to0$. At $s=s_*$ we know that $G^{(2+\varepsilon)}_{s_*^2}(\pi)=0$, hence \eqref{expand} gives
\[
  e^{-\pi s_*}=\frac{\varepsilon A_0}{\pi s_*^3}(1+o(1)).
\]
Taking logarithms first gives $\pi s_*\sim V$, then
$\pi s_*=V+3\log V+\log(1/(\pi^2 A_0))+o(1)$, as claimed.
\end{proof}

\appendix
\section{Zhu's Jacobi--Stirling conjecture}

For $a,b\geq 0$, set $T_{0,0}=1$, set $T_{n,k}=0$ outside
$0\leq k\leq n$, and define
\[
 T_{n+1,k}=T_{n,k-1}+\lambda_k T_{n,k},\qquad
 \lambda_k=ak^2+bk,\qquad M_n(q)=\sum_{k=0}^nT_{n,k}q^k .
\]
In Section 5.2 of his 2024 paper, Zhu records his 2018 conjecture
that every minor of $(M_{i+j}(q))_{i,j\geq0}$ has nonnegative
coefficients in $a,b,q$. Deb's problem list records the same assertion
as P1.9. Zhu's Proposition 5.7 gives a positive result for the row sums
with the extra weight $k!$ \cite{zhu,deb}. Manuscript [167] in the corpus already refutes the unweighted conjecture by the following determinant and gives the general three-node computation \cite{vmrc}. We reproduce this earlier result to specify the starting point of the numerical classification.

\begin{proposition}[A quadratic obstruction]
The first three rows and columns of the unweighted Hankel matrix have determinant
\[
 \det(M_{i+j}(q))_{0\leq i,j\leq2}
 =q^3\{2(a+b)(2a+b)(3a+b)-4a^2q\}.
\]
Consequently Zhu's conjecture is false. For each fixed $a>0$, $b\geq0$,
the same minor is negative for every $q>(a+b)(2a+b)(3a+b)/(2a^2)$.
\end{proposition}
\begin{proof}
For any node sequence starting at zero, write $A=\lambda_1$,
$B=\lambda_2$, $C=\lambda_3$. The recurrence gives
\[
\begin{aligned}
 M_0&=1,\qquad M_1=q,\qquad M_2=Aq+q^2,\\
 M_3&=A^2q+(A+B)q^2+q^3,\\
 M_4&=A^3q+(A^2+AB+B^2)q^2+(A+B+C)q^3+q^4.
\end{aligned}
\]
Eliminating the first row and column gives the determinant
$(M_2-M_1^2)(M_4-M_2^2)-(M_3-M_1M_2)^2$. Substitution yields
\[
 q^3\bigl[AB(B-A)+q\{A(B+C-A)-B^2\}\bigr].
\]
Inserting the quadratic nodes gives the asserted negative coefficient.
At $(a,b,q)=(1,0,4)$ the determinant is exactly $-256$.
\end{proof}

\section{The exact range of positive moment measures}

The determinant leaves many numerical specializations undecided.
The auxiliary problem is to recover a measure supported on the
interpolation nodes from the values of its Newton polynomial moments,
and decide when the candidate measure is positive.
A Stieltjes moment sequence is a sequence of the moments
$\int_0^\infty u^n\,d\mu(u)$ of a positive measure.
For $\nu\geq0$, write $j_{\nu,m}$ for the $m$th strictly positive
zero of the Bessel function $J_\nu$.

\begin{theorem}[Numerical classification]\label{thm:js-range}
For $a>0$, $b\geq0$, and $q\geq0$, the sequence $(M_n(q))_{n\geq0}$ is
a Stieltjes moment sequence if and only if
\[
      0\leq q\leq \frac a4\,j_{b/a,1}^2.
\]
For $a=0$, it is a Stieltjes moment sequence for every $b,q\geq0$.
\end{theorem}

\begin{lemma}[The masses at the quadratic nodes]\label{lem:js-masses}
Assume $a,q>0$. Put $s=b/a$, $r=\sqrt{q/a}$ and
$c_j=(2j+s)\Gamma(j+s)/j!$.
At $j=s=0$ use $c_0=1$ by continuity.
A signed measure with moments $M_n(q)$ is
\[
     \eta=\sum_{j=0}^\infty w_j\delta_{\lambda_j},
     \qquad w_j=c_jr^{-s}J_{2j+s}(2r).
\]
For every $h>0$, $\int e^{h\sqrt{u}}\,d|\eta|(u)<\infty$.
\end{lemma}
\begin{proof}
In the Newton basis $P_k(u)=\prod_{i=0}^{k-1}(u-\lambda_i)$,
the relation $uP_k=P_{k+1}+\lambda_kP_k$ gives
$u^n=\sum_{k=0}^n T_{n,k}P_k(u)$. The Lagrange formula for its
divided differences gives, with $0^0=1$,
\[
 T_{n,k}=\sum_{j=0}^k
   \frac{\lambda_j^n}{\prod_{0\leq i\leq k,\ i\ne j}(\lambda_j-\lambda_i)} .
\]
The right hand side vanishes for $k>n$ by polynomial interpolation.
The explicit finite-sum formula for Jacobi--Stirling entries also
appears in Caki\'c, El-Desouky and Gomaa \cite{cakic}.
Writing $k=j+\ell$ and using
$\lambda_j-\lambda_i=a(j-i)(j+i+s)$, the coefficient contributed
to the mass at $\lambda_j$ by $q^kT_{n,k}$ after removal of $\lambda_j^n$ is
\[
     D_{j,\ell}=
     \frac{(-1)^\ell c_jr^{2j+2\ell}}{\ell!\,\Gamma(2j+s+\ell+1)} .
\]
Because $s\geq0$, the absolute sum over $\ell$ is bounded by
\[
 \sum_{\ell\geq0}|D_{j,\ell}|
 \leq \frac{e^{r^2} c_jr^{2j}}{\Gamma(2j+s+1)}.
\]
The $c_j$ have at most polynomial growth by Stirling's formula.
The last bound is summable over $j$ after multiplication by
$e^{h\sqrt{\lambda_j}}$ or $\lambda_j^n$. This justifies summation
in either order, including all the added zero divided differences.
The Bessel power series \cite[10.2.2]{dlmf} now gives
$\sum_{\ell\geq0}D_{j,\ell}=c_jr^{-s}J_{2j+s}(2r)$.
At $n=0$ its normalization is the classical Neumann expansion
\cite[10.23.15]{dlmf}.
\end{proof}

\begin{lemma}[A sign criterion for discrete moment expansions]
If a finite signed measure $\eta$ on $[0,\infty)$ has
$\int e^{h\sqrt{u}}\,d|\eta|(u)<\infty$ for some $h>0$,
its moments admit a positive measure on $[0,\infty)$ if and only if
$\eta$ itself is positive.
\end{lemma}
\begin{proof}
Write $\eta=\eta^+-\eta^-$. If a positive $\rho$ has the same
moments as $\eta$, then $\rho+\eta^-$ and $\eta^+$ have the same
moments. With $K=\int e^{h\sqrt{u}}\,d\eta^+(u)$, their common
$n$th moment satisfies $m_n\leq K(2n)!h^{-2n}$. Hence
$\sum_{n\geq1}m_n^{-1/(2n)}=\infty$. The Stieltjes form of
Carleman's criterion gives $\rho+\eta^-=\eta^+$.
Zero moments give the same conclusion directly for a measure supported
at zero. Thus $\rho=\eta$. The converse uses $\rho=\eta$.
\end{proof}

\begin{proof}[Proof of Theorem~\ref{thm:js-range}]
For $a,q>0$ the nodes are distinct and $c_j>0$. The lemmas reduce
the question to simultaneous nonnegativity of
$J_{s+2j}(2r)$ for all integers $j\geq0$.
For $\nu\geq0$, the function $J_\nu$ is positive before its
first positive zero, its positive zeros are simple, and
$j_{\nu,1}$ increases strictly with $\nu$ and tends to infinity
\cite[10.21(i),(iv),(vii)]{dlmf}.
It follows that all the required signs are nonnegative for
$0<2r\leq j_{s,1}$.

For the converse, the second zero $z=j_{\nu,2}$ is an upward
crossing. The Bessel recurrences \cite[10.6.1--2]{dlmf} give
\[
 J_{\nu+2}(z)=\frac{2(\nu+1)}{z}J_{\nu+1}(z)
             =-\frac{2(\nu+1)}{z}J_\nu'(z)<0.
\]
Thus $j_{\nu+2,1}<j_{\nu,2}$.
The open intervals $(j_{s+2j,1},j_{s+2j,2})$ overlap successively,
and their left endpoints increase without bound. They cover
$(j_{s,1},\infty)$, and on each interval the corresponding
Bessel function is strictly negative. This proves the cutoff,
including equality at the first zero.

At $q=0$, the representing measure is $\delta_0$. At $a=b=0$,
it is $\delta_q$. If $a=0<b$, the recurrence gives
$M_{n+1}=qM_n+bqM_n'$. Its exponential generating function is
$\exp[(q/b)(e^{bt}-1)]$, which is the moment generating function
of $bN$ for $N$ Poisson of mean $q/b$.
\end{proof}

The negative coefficient was already recorded in the corpus \cite{vmrc}.
The first-zero criterion answers the different question of
which numerical specializations retain a positive measure.
For example, at $a=1,b=0$ the order-three determinant stays
nonnegative up to $q=3$, while the true endpoint is $j_{0,1}^2/4$.
\section*{Scope and priority}
The principal named target is Le--Pelinovsky Conjecture 1.1. The proof uses the original Fourier multiplier and the original parameter quantifiers. The strict circle comparison supplies a true first contact; the sign of all noncentral translates resolves the fate of a multiple zero, which convergence alone leaves open. These are the steps on which the additional conclusions depend. The threshold expansion in Theorem~\ref{threshold-asym} gives error control and the next terms for the formal approximation in Appendix C of \cite{lp}. The two-term growth law concerns the iterated limit in which $c\to\infty$ first and $\alpha\uparrow4$ second.
  
For Conjecture 1.2 the line tail and the sign change were already available in \cite{sk}; the holomorphic representation supplies finite counting and the local classification. The fourth-order boundary case is explicitly credited to \cite{lp}. The Jacobi--Stirling appendix extends the numerical range determined by the corpus. Its original counterexample belongs to manuscript [167]. The exact numerical range is an auxiliary result and carries no claim of settling a second previously open named conjecture.

The literature comparison included the original papers, Simon's explicit open-problem remark, subsequent citing records, searches for the Fourier multiplier, signed resolvent asymptotics, periodic zero counts, and the corpus at the revision below. Rotation of contours and periodization already occur in Simon's work. Exponential oscillations competing with algebraic decay occur in the fractional-oscillator literature \cite{gm}. The local rescaling, lobe counts and splitting of a double zero also have antecedents in the corpus. For ordinary even-order operators, Cabada and López-Somoza study constant sign and antipodal contact \cite{cl}. We found no earlier proof of the periodic stabilization or the combined first-contact and counting theorem in the works examined. Historical priority and external validation remain unconfirmed.
\begin{thebibliography}{99}
\bibitem{lp} U. Le and D. E. Pelinovsky, Green's function for the fractional KdV equation on the periodic domain via Mittag-Leffler function, \emph{Fract. Calc. Appl. Anal.} 24 (2021), 1507--1534. \href{https://doi.org/10.1515/fca-2021-0063}{DOI 10.1515/fca-2021-0063}. \href{https://arxiv.org/pdf/2101.02269}{Author version}, Conjectures 1--2, Appendix B.
\bibitem{simon} T. Simon, On the Green function of the killed fractional Laplacian on the periodic domain, \emph{Fract. Calc. Appl. Anal.} 24 (2021), 1629--1635. \href{https://arxiv.org/abs/2109.01837}{arXiv:2109.01837}.
\bibitem{sk} A. G. Stefanov and P. G. Kevrekidis, Kinks of fractional $\phi^4$ models: existence, uniqueness, monotonicity, stability, and sharp asymptotics. \href{https://arxiv.org/abs/2503.15806}{arXiv:2503.15806} (2025), Lemma 6 and subsequent remarks. \href{https://doi.org/10.1088/1361-6544/ae669c}{Journal version}.
\bibitem{gm} R. Gorenflo and F. Mainardi, Fractional calculus: integral and differential equations of fractional order, in \emph{Fractals and Fractional Calculus in Continuum Mechanics}, Springer (1997), 223--276. \href{https://arxiv.org/abs/0805.3823}{Author version}.
\bibitem{cl} A. Cabada and L. López-Somoza, Spectral characterization of the constant sign Green's functions for periodic and Neumann boundary value problems of even order, \emph{Differ. Equ. Appl.} 14 (2022), 335--347. \href{https://doi.org/10.7153/dea-2022-14-24}{DOI 10.7153/dea-2022-14-24}.
\bibitem{vmrc} QEEDDD, \emph{Verifiable Mathematical Reasoning Corpus}, \href{https://github.com/QEEDDD/Verifiable-Mathematical-Reasoning-Corpus/tree/5887cdea4898f0a3aa5bf7a8e0829697bd2b0f4b}{revision 5887cdea4898}. See [026] on rescaled determinant zeros, [123] on a double zero and its splitting, [127] on Mittag-Leffler lobe counts, [167] on the Jacobi--Stirling determinant, and [175] on Mittag-Leffler estimates on a half-line.
\bibitem{zhu} B.-X. Zhu, Coefficientwise Hankel-total positivity of the row-generating polynomials for the output matrices of certain production matrices, \emph{Adv. Math.} 447 (2024), 109695, Section 5.2. \href{https://arxiv.org/html/2202.03793v2}{Author version}.
\bibitem{deb} B. Deb, \emph{Positivity problems in enumerative and algebraic combinatorics}, \href{https://bishaldeb.com/positivity/themes/SMSCHTP/hankel.html}{Problem P1.9}.
\bibitem{cakic} N. P. Caki\'c, B. S. El-Desouky and R. S. Gomaa, New family of Jacobi-Stirling numbers, \emph{Appl. Anal. Discrete Math.} 17 (2023), 538--547. \href{https://doi.org/10.2298/AADM210829013C}{DOI 10.2298/AADM210829013C}.
\bibitem{dlmf} NIST, \emph{Digital Library of Mathematical Functions}, \href{https://dlmf.nist.gov/10}{Chapter 10}, especially Sections 10.2, 10.6, 10.21 and 10.23.
\end{thebibliography}
\end{document}
